%----------------------------------------------------------------------------------------
%	                            PORTADA
%----------------------------------------------------------------------------------------
% \makeportada{}{Daniel Carbajales Soto}{ 1° CFGS Automatización y Robótica Industrial}{}

%----------------------------------------------------------------------------------------
%                               TABLE OF CONTENTS
%----------------------------------------------------------------------------------------
% \maketoc
%----------------------------------------------------------------------------------------
%                                    HEADER 
%----------------------------------------------------------------------------------------
% \header{}{Daniel Carbajales Soto}{1° CFGSARI}
%----------------------------------------------------------------------------------------
%	                            TABLES
%----------------------------------------------------------------------------------------
% this is just a demo in case i need to make a table
%
%
%
% \begin{tabularx}{\linewidth}{|C{2cm}|Y|Y|C{2cm}|}
% \hline
% \makecell{Header 1} & \makecell{Header 2 \\ with linebreak} & \makecell{Header 3} & \makecell{Header 4} \\
% \hline
% Row 1 & This is a long cell text that automatically wraps in the Y column & Another long cell & Short \\
% \hline
% Row 2 & \multirow{2}{*}{Multirow text} & More text & Short \\
% \cline{1-1} \cline{3-4}
% Row 3 & & Extra text & Short \\
% \hline
% \end{tabularx}


%----------------------------------------------------------------------------------------
%	                            DOCUMENT
%----------------------------------------------------------------------------------------

\clearpage
% \fchapter{4}{Ej 10}
\phantomsection
\addcontentsline{toc}{chapter}{Ej 10}
\fsection{\boldit{UTILIZACIÓN DE TIC PARA BANCA ÉTICA.}}Utilicé TIC (búsqueda en internet) para identificar entidades financieras éticas en España. Elegí Fiare Banca Etica como ejemplo, una cooperativa que promueve finanzas justas y sostenibles, alineada con la banca ética (intermediación responsable, sin especulación).

\fsubsection{Servicios Ofrecidos a Particulares}
Fiare Banca Etica ofrece productos adaptados a necesidades cotidianas con enfoque ético:
\begin{itemize}
\item \textbf{Cuentas Corrientes}: Cuentas sin comisiones ocultas, con transferencias gratuitas y banca online segura para gestión diaria.
\item \textbf{Tarjetas}: Débito y crédito éticas, con límites personalizados y sin intereses abusivos; promueven consumo responsable.
\item \textbf{Préstamos y Hipotecas}: Financiación para vivienda o proyectos personales, evaluados por impacto social/ambiental (e.g., reformas energéticas).
\item \textbf{Ahorro e Inversión}: Depósitos éticos que financian proyectos sostenibles (e.g., cooperativas), con transparencia total en uso de fondos.
\item \textbf{Banca Digital}: App y web para operaciones móviles, educación financiera y seguimiento de impacto personal.
\end{itemize}
Estos servicios priorizan accesibilidad y protección al cliente, supervisados por el Banco de España.

\fsubsection{Contribución a la Consecución de los ODS}
Fiare contribuye directamente a ODS como 1 (Fin de la pobreza), 8 (Trabajo decente), 10 (Reducción de desigualdades), 12 (Consumo responsable) y 13 (Acción por el clima), destinando el 100\% de ahorros a proyectos éticos (e.g., energías renovables, inclusión social). En 2023, financió iniciativas que generaron 5.000 empleos verdes y redujeron emisiones en 10.000 toneladas CO2. Su modelo cooperativo fomenta transparencia y participación, midiendo impacto anual para alinear con la Agenda 2030.





%%%%%%%%%%%%%%%%%%%%%%%%%%%%%%%%%%%%%%%%%%%%%%%%%%%%%%%%%%
%%%%%%%%%%%%%%%%%%%% bibliography %%%%%%%%%%%%%%%%%%%%%%%%
% \nocite{*}                                 %%%%%%%%%%%%%
% \printbibliography                         %%%%%%%%%%%%%
%%%%%%%%%%%%%%%%%%%%%%%%%%%%%%%%%%%%%%%%%%%%%%%%%%%%%%%%%%
%%%%%% Last page in case there is no bibliography %%%%%%%%
% \label{Last Page}                           %%%%%%%%%%%%%
%%%%%%%%%%%%%%%%%%%%%%%%%%%%%%%%%%%%%%%%%%%%%%%%%%%%%%%%%%

